\section{The entire equality locus and ordinary returns}
\label{sec:equality}

\begin{proof}[Proof of the equality assertion in Theorem~\ref{thm:main}]
The affine-root and large-diameter cases have at most nine points.
The cases using only two symbols have at most four.
Thus a map with eleven periodic points must have its global
normalization at the unique tuple \eqref{eq:max-tuple}.
The normalized polynomial is
\[
 g(t)=6-t+t(t-4)(t-2)
     =t^3-6t^2+7t+6
     =(t-2)^3-5(t-2)+4.
\]
Undoing $g(t)=f(t+m)-2m$ and putting $h=m+2$ gives precisely
\eqref{eq:equality-family}. In the original coefficients it is
\begin{equation}
 b=-3h,\qquad c=3h^2-5,\qquad a=-h^3+7h.
 \label{eq:equality-coefficients}
\end{equation}
This proves necessity for every original coefficient triple.

Conversely, for $f_0(t)=t^3-5t$ the disjoint cycles are
\begin{equation}
 (0),\quad(-2,1),\quad(-1,2),\quad
 (-2,0,2),\quad(-2,2,0).
 \label{eq:eleven-cycles}
\end{equation}
Substitution verifies them, and their lengths $1,2,2,3,3$ give
eleven distinct points. The proved upper bound excludes any
additional rational periodic points. Conjugating by $T_h$
translates every coordinate in \eqref{eq:eleven-cycles} by $h$
and gives exactly \eqref{eq:equality-family}.
Hence every integer $h$ realizes equality, with no omitted
exceptional coefficient branch.
\end{proof}

The ordinary return series follows from the cycle classification;
it is not a separate arithmetic counting problem. Let $c_j(f)$ be
the number of distinct cycles of least period $j$, where
$j\in\{1,2,3,4,6\}$, as obtained from Table~\ref{tab:templates}.
Cycles related only by reversal remain distinct when reversal is
not already a rotation. Define
\[
 N_n(f)=\#\Fix(H_f^n,\Q^2),\qquad
 Z_f(z)=\exp\left(\sum_{n\ge1}\frac{N_n(f)}{n}z^n\right).
\]
These counts use actual rational points, each with weight one.
Every point fixed by an iterate is periodic, so all counts are finite.

\begin{corollary}[Native ordinary zeta]
\label{cor:zeta}
As identities of formal power series,
\[
 N_n(f)=\sum_{\substack{j\in\{1,2,3,4,6\}\\j\mid n}}j\,c_j(f),
 \qquad
 Z_f(z)=\prod_{j\in\{1,2,3,4,6\}}(1-z^j)^{-c_j(f)}.
\]
On the eleven-point equality family,
\[
 Z_f(z)=\frac{1}{(1-z)(1-z^2)^2(1-z^3)^2}.
\]
\end{corollary}
\begin{proof}
A cycle of length $j$ contributes all of its $j$ points to the
fixed set of the $n$th iterate exactly when $j\mid n$.
This proves the first identity. Its contribution to the logarithm
of the zeta is
\[
 \sum_{m\ge1}\frac{j}{jm}z^{jm}
 =\sum_{m\ge1}\frac{z^{jm}}m=-\log(1-z^j).
\]
There are finitely many cycles, so summing and exponentiating gives
the product. Formula \eqref{eq:eleven-cycles} gives the last identity.
\end{proof}

No scheme multiplicity or auxiliary return clock enters this
corollary. The map is invertible, so a rational point that is
eventually periodic is already periodic: applying the inverse
of an initial iterate removes any purported preperiod.
